\chapter{結論}
\label{chap:conclusion}

\section{結論と貢献}
\label{sec:conclusion}

\draftnote{context.md §6．三つの主結果を一文で束ねる．貢献（足す側から構成的に示した），「この研究から言えないこと」の列挙をここでも省かない——頑健性（\ref{sec:robustness}節）の結果が出たら，束ねる一文はそれを含む形に書き直す．}

\section{限界}
\label{sec:limitations}

\draftnote{context.md §7．設計の一点性（$\lambda=0.1$・第16層・単一モデル・単一データセット・1エポック），署名が解き方の正しい形式化かは未検証（GQA 公式 consistency との食い違い52.1\%の読み方を含む），「使われていない」は第16層・一次近似の範囲であること．}

\section{今後の課題}
\label{sec:future}

\draftnote{context.md §8．執筆時点で残っているものだけを書く（済んだ測定は\ref{chap:results}章へ移す）：稀な署名の表現側の層別（§8.2），答えの型ラベルでの整合——主語の確定（§8.3），$\lambda$ と層を振る・別モデルと別データセット（§8.4，署名に代わる形式化が要ること）．}
